\documentclass{beamer}

\usepackage{amsmath,amsfonts,amssymb}
\usepackage{physics}
\usepackage{graphicx}
\usepackage{bm}

\title{Bootstrap Resampling Methods}
\subtitle{Statistical Error Estimation in Computational Physics}
\author{Morten Hjorth-Jensen}
\date{Spring 2026}

\begin{document}

\frame{\titlepage}

%------------------------------------------------

\begin{frame}{Motivation}

In many areas of physics we estimate observables from finite datasets:

\[
X_1, X_2, \dots, X_N
\]

Examples:

\begin{itemize}
\item Monte Carlo simulations
\item Lattice gauge theory
\item Molecular dynamics
\item Experimental measurements
\end{itemize}

Goal: estimate statistical uncertainty of an estimator

\[
\hat{\theta} = f(X_1,\dots,X_N)
\]

Problem:

\begin{itemize}
\item Analytical variance formulas may be unavailable
\item Estimators may be nonlinear
\item Distribution of $\hat{\theta}$ may be unknown
\end{itemize}

Bootstrap resampling provides a numerical solution.
\end{frame}

%------------------------------------------------

\begin{frame}{Basic Idea of the Bootstrap}

Suppose the unknown distribution generating the data is

\[
F(x)
\]

but we only observe samples

\[
X_1,\dots,X_N \sim F
\]

The bootstrap replaces the unknown distribution by the empirical distribution

\[
\hat{F}_N(x)
\]

defined by

\[
\hat{F}_N(x)=\frac{1}{N}\sum_{i=1}^{N}\delta(x-X_i)
\]

Interpretation:

\begin{itemize}
\item Treat the observed dataset as the population
\item Generate new samples by resampling with replacement
\end{itemize}

This produces synthetic datasets from $\hat{F}_N$.
\end{frame}

%------------------------------------------------

\begin{frame}{Bootstrap Resampling}

Generate a bootstrap sample

\[
X_1^*, X_2^*, \dots, X_N^*
\]

by drawing $N$ values from the original data

\[
\{X_1,\dots,X_N\}
\]

with replacement.

Example:

\[
\{1.2, 0.8, 1.4, 0.9\}
\]

Possible bootstrap sample:

\[
\{0.8,1.4,1.4,1.2\}
\]

Each bootstrap dataset mimics an independent experiment.
\end{frame}

%------------------------------------------------

\begin{frame}{Bootstrap Distribution of an Estimator}

For each bootstrap sample compute

\[
\hat{\theta}^* = f(X_1^*,\dots,X_N^*)
\]

Repeat this procedure $B$ times:

\[
\hat{\theta}_1^*, \hat{\theta}_2^*, \dots, \hat{\theta}_B^*
\]

These values form the \textbf{bootstrap distribution}

\[
P_B(\hat{\theta})
\]

which approximates the sampling distribution of the estimator

\[
P(\hat{\theta})
\]

in the limit of large $N$.
\end{frame}

%------------------------------------------------

\begin{frame}{Bootstrap Estimate of the Variance}

The bootstrap estimate of the variance is

\[
\mathrm{Var}_{\text{boot}}(\hat{\theta})
=
\frac{1}{B-1}
\sum_{b=1}^{B}
(\hat{\theta}_b^*-\bar{\theta}^*)^2
\]

where

\[
\bar{\theta}^*=
\frac{1}{B}\sum_{b=1}^{B}\hat{\theta}_b^*
\]

Thus the bootstrap standard error becomes

\[
\sigma_{\text{boot}}
=
\sqrt{\mathrm{Var}_{\text{boot}}}.
\]

As $B\to\infty$, the estimate converges to the true variance.
\end{frame}

%------------------------------------------------

\begin{frame}{Consistency of the Bootstrap}

Let

\[
\hat{\theta}=f(X_1,\dots,X_N)
\]

be an estimator.

Under mild regularity conditions,

\[
\sqrt{N}(\hat{\theta}-\theta)
\]

has a limiting distribution.

Bootstrap theory shows

\[
\sqrt{N}(\hat{\theta}^*-\hat{\theta})
\]

approximates the same limiting distribution.

Thus

\[
P^*(\hat{\theta}^*-\hat{\theta})
\approx
P(\hat{\theta}-\theta).
\]

This is the theoretical justification of the bootstrap.
\end{frame}

%------------------------------------------------

\begin{frame}{Bootstrap Confidence Intervals}

Let

\[
\hat{\theta}_1^*,\dots,\hat{\theta}_B^*
\]

be sorted.

\textbf{Percentile interval}

For confidence level $1-\alpha$:

\[
[\hat{\theta}_{\alpha/2}^*, \hat{\theta}_{1-\alpha/2}^*]
\]

Example:

95\% interval

\[
[\hat{\theta}_{2.5\%}^*,\hat{\theta}_{97.5\%}^*]
\]

This method uses the empirical quantiles of the bootstrap distribution.
\end{frame}

%------------------------------------------------

\begin{frame}{Bias Estimation}

Bootstrap also estimates bias:

\[
\text{Bias} =
\mathbb{E}[\hat{\theta}] - \theta
\]

Bootstrap estimate

\[
\widehat{\text{Bias}}
=
\bar{\theta}^* - \hat{\theta}
\]

Bias-corrected estimator:

\[
\hat{\theta}_{BC}
=
\hat{\theta} - \widehat{\text{Bias}}
\]

Bootstrap therefore provides both

\begin{itemize}
\item variance estimates
\item bias corrections
\end{itemize}
\end{frame}

%------------------------------------------------

\begin{frame}{Bootstrap Algorithm}

Given data

\[
X_1,\dots,X_N
\]

\begin{enumerate}

\item Choose number of resamples $B$

\item For $b=1\dots B$

\begin{itemize}

\item Draw bootstrap sample with replacement
\[
X_1^*,\dots,X_N^*
\]

\item Compute estimator
\[
\hat{\theta}_b^*
\]

\end{itemize}

\item Compute

\[
\sigma_{\text{boot}}
\]

\end{enumerate}

Typical choice:

\[
B = 1000 - 10000
\]
\end{frame}

%------------------------------------------------

\begin{frame}{Bootstrap Pseudocode}

\begin{verbatim}
theta_boot = []

for b in range(B):

    sample = random.choice(data, size=N)

    theta_star = estimator(sample)

    theta_boot.append(theta_star)

mean_boot = mean(theta_boot)

variance_boot = var(theta_boot)
\end{verbatim}

The list \texttt{theta\_boot} approximates the sampling distribution of the estimator.
\end{frame}

%------------------------------------------------

\begin{frame}{Applications in Physics}

Bootstrap is widely used in

\begin{itemize}

\item Lattice QCD
\item Quantum Monte Carlo
\item Molecular simulations
\item Experimental particle physics
\item Parameter estimation in model fitting

\end{itemize}

Examples of estimators:

\begin{itemize}

\item correlation functions
\item energy differences
\item nonlinear observables
\item fitted parameters
\end{itemize}
\end{frame}

%------------------------------------------------

\begin{frame}{Example: Monte Carlo Observable}

Suppose we estimate an observable

\[
\langle O \rangle
=
\frac{1}{N}\sum_i O_i
\]

But we want uncertainty for a nonlinear function

\[
R = \frac{\langle A \rangle}{\langle B \rangle}
\]

Analytic variance formulas become complicated.

Bootstrap procedure:

\begin{itemize}

\item resample pairs $(A_i,B_i)$
\item recompute ratio
\item obtain distribution of $R$

\end{itemize}

This yields robust uncertainty estimates.
\end{frame}

%------------------------------------------------

\begin{frame}{Bootstrap vs Other Resampling Methods}

Common resampling techniques:

\begin{itemize}

\item Bootstrap
\item Jackknife
\item Blocking
\end{itemize}

Bootstrap advantages:

\begin{itemize}

\item works for complex estimators
\item easy to implement
\item provides full distribution
\end{itemize}

Bootstrap limitation:

\begin{itemize}

\item assumes independent samples
\end{itemize}

For correlated Monte Carlo data we combine:

\begin{itemize}

\item blocking + bootstrap
\end{itemize}
\end{frame}

%------------------------------------------------

\begin{frame}{Summary}

Bootstrap resampling is a powerful statistical tool.

\begin{itemize}

\item Approximates sampling distributions numerically
\item Requires minimal analytical assumptions
\item Provides variance, bias, and confidence intervals
\item Widely used in computational physics

\end{itemize}

In modern simulations bootstrap analysis is a standard component of data analysis pipelines.
\end{frame}

%------------------------------------------------

\end{document}
